\documentclass[10pt,twocolumn]{article}

\usepackage[letterpaper,top=0.67in,bottom=0.72in,left=0.67in,right=0.67in,columnsep=0.24in]{geometry}
\usepackage[T1]{fontenc}
\usepackage[utf8]{inputenc}
\usepackage{lmodern}
\usepackage{microtype}
\usepackage{graphicx}
\usepackage{balance}
\usepackage{float}
\usepackage{placeins}
\usepackage{booktabs}
\usepackage{longtable}
\usepackage{tabularx}
\usepackage{array}
\usepackage{ragged2e}
\usepackage{etoolbox}
\usepackage{amsmath,amssymb}
\usepackage[numbers,sort&compress]{natbib}
\usepackage[font=small,labelfont=bf,labelsep=period]{caption}
\usepackage{xcolor}
\usepackage{titlesec}
\usepackage{fancyhdr}
\usepackage[hidelinks,breaklinks=true]{hyperref}
\usepackage[nameinlink,noabbrev]{cleveref}

\definecolor{chemblue}{HTML}{26577C}
\definecolor{chemlink}{HTML}{234F70}
\hypersetup{
  unicode=true,
  pdfencoding=auto,
  colorlinks=true,
  linkcolor=chemlink,
  citecolor=chemlink,
  urlcolor=chemlink,
  pdftitle={Dissecting Executable Scientific Intelligence with
Controlled World-Model Interventions},
  pdfauthor={Jiangjie Qiu; Yijun Li; Yaotian Yang; Honghao Chen; Wentao
Li; Xiaonan Wang},
  pdfsubject={Causal dissection of experimental search, predictive
correction, executable laws and action transfer in AI agents},
  pdfkeywords={AI scientist; autonomous experimentation; initial world
model; scientific priors; law discovery; counterfactual prediction;
action transfer; chemical worlds}
}

\setlength{\parindent}{1em}
\setlength{\parskip}{0pt}
\setlength{\emergencystretch}{1.5em}
\raggedbottom
\setlength{\abovedisplayskip}{6pt}
\setlength{\belowdisplayskip}{9pt}
\setlength{\abovedisplayshortskip}{4pt}
\setlength{\belowdisplayshortskip}{9pt}
\setlength{\textfloatsep}{8pt plus 2pt minus 2pt}
\setlength{\dbltextfloatsep}{8pt plus 2pt minus 2pt}
\setlength{\floatsep}{7pt plus 2pt minus 2pt}
\setlength{\dblfloatsep}{7pt plus 2pt minus 2pt}
\setlength{\intextsep}{7pt plus 2pt minus 2pt}
\setcounter{topnumber}{3}
\setcounter{dbltopnumber}{3}
\setcounter{totalnumber}{5}
\renewcommand{\topfraction}{0.94}
\renewcommand{\dbltopfraction}{0.94}
\renewcommand{\textfraction}{0.06}
\renewcommand{\floatpagefraction}{0.82}
\renewcommand{\dblfloatpagefraction}{0.82}
\captionsetup{justification=justified,singlelinecheck=false}
\captionsetup[figure]{skip=5pt}
\captionsetup[table]{skip=4pt}
\newcolumntype{L}[1]{>{\RaggedRight\arraybackslash}p{#1}}
\newcolumntype{Y}{>{\RaggedRight\arraybackslash}X}
\clubpenalty=8000
\widowpenalty=8000
\displaywidowpenalty=8000
\titleformat{\section}{\large\bfseries\color{chemblue}\raggedright\hyphenpenalty=10000\exhyphenpenalty=10000}{}{0pt}{}
\titleformat{\subsection}{\normalsize\bfseries\raggedright\hyphenpenalty=10000\exhyphenpenalty=10000}{}{0pt}{}
\titlespacing*{\section}{0pt}{10pt plus 2pt minus 1pt}{4pt}
\titlespacing*{\subsection}{0pt}{7pt plus 2pt minus 1pt}{3pt}

\pagestyle{fancy}
\fancyhf{}
\fancyhead[L]{\footnotesize Experimental Intelligence in AI Scientists}
\fancyhead[R]{\footnotesize Working manuscript}
\fancyfoot[C]{\thepage}
\renewcommand{\headrulewidth}{0.25pt}

\providecommand{\tightlist}{%
  \setlength{\itemsep}{0pt}\setlength{\parskip}{0pt}}
\AtBeginEnvironment{thebibliography}{\small}

\title{\vspace{-1.8em}\textbf{Dissecting Executable Scientific
Intelligence\\with Controlled World-Model Interventions}}
\author{%
Jiangjie Qiu\textsuperscript{1,\textdagger}, Yijun
Li\textsuperscript{1,\textdagger}, Yaotian
Yang\textsuperscript{1,\textdagger}, Honghao
Chen\textsuperscript{1}, Wentao Li\textsuperscript{1}, Xiaonan
Wang\textsuperscript{1,*}%
\\[0.45em]
\parbox{0.92\textwidth}{\centering\small
\textsuperscript{1}Beijing Key Laboratory of Artificial Intelligence for
Advanced Chemical Engineering Materials, State Key Laboratory of
Chemical Engineering and Low-Carbon Technology, Department of Chemical
Engineering, Tsinghua University, Beijing 100084, China}%
\\[0.22em]\small Jiangjie Qiu, Yijun Li and Yaotian Yang contributed
equally.%
\\[0.35em]\small *Correspondence: \texttt{wangxiaonan@tsinghua.edu.cn}%
}
\date{}

\begin{document}

\twocolumn[
\begin{@twocolumnfalse}
\maketitle
\begin{abstract}
Scientific agents enter experiments with an initial world model that can
guide useful search but can also be wrong. Endpoint success alone cannot
distinguish correct scientific inference from a favorable heuristic
trajectory. We organized a completed three-part programme around three
questions: do priors change search and endpoints; does evidence
selectively repair wrong priors and yield executable regularities; and
can learned information support an unseen action? In a prospective
cohort of 135 sessions and 1,243/1,260 planned experiments, prior
effects were strongly task dependent. Correct entity information
produced a durable advantage in partition, whereas correct structural
information gave crystallization an initial head start that narrowed
during exploration. Structural partition showed that explicit
organization and correctness can contribute separately. Prediction error
generally decreased, but selective correction failed at the entity
(\(p=0.990\)), parametric (\(p=0.079\)) and structural (\(p=1.000\))
loci. In matched-evidence studies, the structural prediction-update
contrast was positive but mixed (+0.0645; 3/5 worlds), while 0/5
summaries recovered the prespecified power law. Executable laws were
better than final explicit predictions in only 50/135 cells and worse in
84/135. A formal multi-task open-action assay completed 45/45 scheduled
cell records and 240/240 truth plus exact-replay queries; among 42
eligible cells, 11 selected the true top-ranked plan, with 0/42
combining an adequate law and a correct action. These results separate
evidence acquisition, numerical revision, structural-law identification,
law compression and unseen-action selection. Private replication and
context-reset artifact portability are future studies, not results of
the present paper.
\end{abstract}
\vspace{0.6em}
\end{@twocolumnfalse}
]

\section{1. Introduction}\label{introduction}

Scientific discovery is not simply the production of a high-scoring
outcome. A useful result can arise from a correct initial model, a
productive but incorrect heuristic, evidence-driven revision or a
fortunate endpoint without a reusable account of the relation that
generated it. These explanations are scientifically distinct even when
the final score is identical. The object of study is therefore the
agent's initial model of what exists, how it works, which quantities
matter and where a learned relation should apply.

AI systems can now plan experiments, call chemistry tools and
participate in closed-loop discovery workflows
\citep{boiko2023autonomous, bran2024augmenting, szymanski2023alab, darvish2025organa, song2025chemagents, vriza2026instruments}.
Interactive scientific environments likewise test hypothesis formation,
intervention and inference
\citep{jansen2024discoveryworld, gandhi2025boxinggym, duan2025scigym, zheng2026newtonbench, yang2026causalab, batzoglou2026replayscm}.
Yet endpoint optimization, belief statements, evidence selection, law
recovery and action transfer are usually entangled.

Executable chemical worlds provide a controlled instrument for
separating these transitions. Entities, process structure, parameters,
measurement mappings and private laws can be instantiated while the
public task, operations, resource opportunities and external world
remain fixed \citep{qiu2026chemworld}. The present study uses that
programmability to intervene on one layer of an initial world model at a
time and to score the resulting transitions. The platform qualification
is not treated as evidence about agent intelligence.

The paper is organized around three questions.

\begin{enumerate}
\def\labelenumi{\arabic{enumi}.}
\tightlist
\item
  \textbf{Do priors change search and endpoints?} We compare opaque,
  aligned and misspecified initial models within matched task--world
  clusters and measure both the first action and the best outcome.
\item
  \textbf{Does evidence selectively repair wrong priors and yield
  executable regularities?} We score held-out predictions,
  matched-evidence updates and executable law summaries separately from
  self-reported confidence.
\item
  \textbf{Can learned information support an unseen action?} We
  distinguish replay of an observed incumbent from ranking fully
  specified candidate plans whose outcomes were hidden during learning.
\end{enumerate}

The completed evidence comprises a prospective cohort, a valid
matched-evidence follow-up and a formal open-action assay. Early
exploratory cohorts, a one-world parametric pilot, an incomplete
prototype, repair trajectories and operational accounting are retained
in the Supplementary Information. Private replication and context-reset
artifact portability are future studies.

\begin{figure*}[!t]
\centering
\includegraphics[width=\textwidth]{figures/prior-discovery/figure-1-prior-to-law.pdf}
\caption{\textbf{The capability chain tested in this paper.} \textbf{a,} A fixed executable world is paired with opaque, aligned or misspecified information at one declared layer. \textbf{b,} A persistent agent selects experiments, observes outcomes and updates predictions and a law summary. \textbf{c,} Participant trajectories and evaluator-owned held-out truth remain separate until scoring. \textbf{d,} Search, prediction, law recovery and unseen-action selection are distinct outcomes; context-reset portability is a future test.}
\label{fig:capability-chain}
\end{figure*}

\section{2. Related work}\label{related-work}

Autonomous chemistry systems demonstrate planning, synthesis, materials
search and closed-loop experimentation on real equipment
\citep{boiko2023autonomous, bran2024augmenting, szymanski2023alab, darvish2025organa, song2025chemagents}.
Virtual laboratories and process environments make experimentation
repeatable and inexpensive
\citep{felton2021summit, hase2021olympus, bloor2024pcgym, beeler2024chemgymrl, malik2026made}.
Interactive discovery benchmarks emphasize repeated cycles of
hypothesis, intervention and inference
\citep{jansen2024discoveryworld, gandhi2025boxinggym, duan2025scigym, nagele2026sciexplorer}.

Prior-availability studies and counterfactual-world benchmarks make
important contributions, but they often change task information without
separating correctness, evidence acquisition and transfer.
Mechanism-oriented systems instead optimize model discovery or
experiment selection
\citep{kabra2026llmautoscilab, murphy2026mda, wahl2026probabilistic}.
Our design is complementary: the external world and action space are
fixed within matched clusters, while one declared component of the
initial scientific model changes. The resulting estimand is a conversion
loss along a capability chain, not a leaderboard score.

\section{3. Conceptual framework}\label{conceptual-framework}

Each matched comparison contains one evaluator-owned world
\(W=(\mathcal{E},G,\Theta,O,C)\) and one participant-facing initial
model
\(M_0=(\widehat{\mathcal{E}},\widehat{G},\widehat{\Theta},\widehat{O},\widehat{S})\).
The public task, action space, observation channels, resource card,
safety rules and stochastic identity are fixed. Only one declared
component of \(M_0\) changes. The manipulated layer can be
entity/ontology, parametric/dynamical or structural/mechanistic;
observation and scope are reserved for separate boundary studies.

Within a layer, information is either opaque, aligned with the fixed
world or equally detailed but misspecified. A misspecified model is a
plausible scientific representation, not a trick question about
obedience. Experimental evidence remains authoritative in all
conditions.

We distinguish five outcomes:

\begin{enumerate}
\def\labelenumi{\arabic{enumi}.}
\tightlist
\item
  endpoint optimization;
\item
  predictive recovery on held-out counterfactual queries;
\item
  selective correction of a wrong prior;
\item
  recovery of an executable law that preserves conditional prediction
  quality; and
\item
  action transfer to a previously unseen, fully specified plan.
\end{enumerate}

The analysis reports transition losses rather than a composite
intelligence score. Endpoint success without predictive and transfer
validity is classified as local optimization rather than law discovery.

\section{4. Completed study
architecture}\label{completed-study-architecture}

The prospective cohort spans five task families with five public worlds
per task for the entity layer, and two validated task families with five
worlds per task for the parametric and structural layers. Every matched
cluster contains opaque, aligned and misspecified arms. Exploratory and
prospective world instances are disjoint. Observation-model
interventions and scope/compositional studies are kept outside the
current denominator.

The completed programme has three evidence blocks:

\begin{enumerate}
\def\labelenumi{\arabic{enumi}.}
\tightlist
\item
  \textbf{Prospective search cohort.} Campaigns contain eight, ten or
  twelve experiments depending on the intervention layer, with
  checkpoints fixed before execution. Each checkpoint records
  predictions, uncertainty, evidence references, a typed law summary and
  the next experimental intent.
\item
  \textbf{Matched-evidence follow-up.} The same contradictory evidence
  is presented to each arm after the initial search. The valid analysis
  combines an unaffected parametric block and a corrected structural
  phase-process block; a superseded run is retained only in the
  Supplementary Information because its evaluator omitted the declared
  intervention.
\item
  \textbf{Formal open-action assay.} After twelve experiments, the agent
  ranks eight fully specified candidate plans with outcomes hidden.
  Public, truth-evaluated and executed plans are verified as identical.
  This assay addresses unseen-action transfer rather than replay of an
  observed incumbent.
\end{enumerate}

\begin{figure*}[!t]
\centering
\includegraphics[width=\textwidth]{figures/prior-discovery/figure-2-formal-cohort.pdf}
\caption{\textbf{Completed evidence architecture.} \textbf{a,} The entity backbone contains five task families and five public worlds per task. \textbf{b,} Each matched task--world cluster contains opaque, aligned and misspecified initial models for one declared locus. \textbf{c,} Parametric and structural blocks are validated separately. \textbf{d,} Prospective search, matched evidence, evaluator scoring and open-action tests retain separate sessions and denominators; private replication and context-reset portability are future work.}
\label{fig:study-architecture}
\end{figure*}

\section{5. Results I --- Priors change search and
endpoints}\label{results-i-priors-change-search-and-endpoints}

The prospective cohort completed \textbf{135/135} scheduled sessions and
\textbf{1,243/1,260} planned experiments. The denominator contains 1,269
closed batch lifecycles, 1,243 final assays and 26 discarded lifecycles.
Failed cells and uncommitted attempts remain in their assigned
denominators.

Prior effects were task dependent. In entity-level partition, aligned
information improved the first experiment by \textbf{0.106} and the best
observed endpoint by \textbf{0.200} relative to misspecified
information, with the same direction in \textbf{5/5} worlds. In
structural crystallization, aligned information improved the first
experiment by \textbf{0.141} in \textbf{5/5} worlds, but the
best-endpoint gap narrowed to \textbf{0.055} and was positive in only
\textbf{3/5} worlds. Structural partition separated organization from
correctness: aligned and misspecified models exceeded opaque identifiers
by \textbf{0.163} and \textbf{0.143}, respectively, while their mutual
difference was only \textbf{0.020}.

The remaining task--locus combinations bounded these cases. Entity-level
reaction safety showed a small aligned-minus-misspecified difference of
\textbf{0.036} at the best endpoint, distillation narrowed from
\textbf{0.053} on the first experiment to \textbf{0.011} at the best
endpoint, and electrochemical and crystallization effects were
heterogeneous. Correct information was therefore neither a universal
performance intervention nor a necessary condition for useful search.

The agents did perform substantive within-session search.
\textbf{91.2\%} of completed experiments used a unique recipe,
\textbf{84.4\%} of session optima occurred after the campaign midpoint
and \textbf{32.6\%} occurred in the final completed experiment. These
observations support search adaptation, not scientific correction: a
productive trajectory can arise while the initial model remains wrong.

\section{6. Results II --- Evidence, selective correction and executable
laws}\label{results-ii-evidence-selective-correction-and-executable-laws}

All sessions submitted five belief checkpoints, yielding
\textbf{675/675} typed records and \textbf{6,300} prespecified
counterfactual predictions. The evaluator completed \textbf{420/420}
truth executions and scored every checkpoint. Prediction error generally
decreased, but the registered selective-correction criterion failed at
all three loci: entity \textbf{\(p=0.990\)}, parametric
\textbf{\(p=0.079\)} and structural \textbf{\(p=1.000\)}. General
prediction improvement and targeted repair of an incorrect prior were
therefore not interchangeable outcomes.

Matched evidence localized the loss without reducing it to a binary
label. In the parametric block, all five misspecified summaries rejected
the supplied high-potential direction after receiving the same
contradictory evidence. In the structural block, the
misspecified-minus-aligned prediction-update contrast was
\textbf{+0.0645}, positive in \textbf{3/5} worlds (exact one-sided
sign-flip \textbf{\(p=0.125\)}), while \textbf{0/5} misspecified
summaries recovered the prespecified power law. Evidence acquisition,
numerical revision and structural-law identification therefore remain
separable transitions.

The evaluator executed all \textbf{135/135} final typed laws. Laws were
more accurate than final explicit predictions in \textbf{50/135} cells,
equal in one and worse in \textbf{84/135}; mean law-minus-final error
was \textbf{+0.069}. Syntax and execution were reliable, but compression
into a reusable relation was often lossy. Blind replay of the committed
recommendation versus the observed incumbent was better, equivalent and
worse in \textbf{1/119/1} of \textbf{121} evaluable cells.
Reproducibility of an observed action was therefore almost complete,
while improvement beyond the observed history was rare.

\begin{figure*}[!t]
\centering
\includegraphics[width=\textwidth]{figures/prior-discovery/figure-3-capability-chain.pdf}
\caption{\textbf{Prediction, law and action dissociation.} Held-out prediction error, selective-correction contrasts, executable-law error and blind incumbent replay are separate readouts. The completed cohort shows broad prediction improvement but no passed selective-correction gate, frequent lossy law compression and almost no recommendation gain over the observed incumbent.}
\label{fig:prediction-law-action}
\end{figure*}

\section{7. Results III --- Learned information and unseen
action}\label{results-iii-learned-information-and-unseen-action}

The formal open-action assay completed \textbf{45/45} scheduled cell
records. Independent evaluation completed \textbf{240/240} truth
executions and \textbf{240/240} exact replays, and verified identity
among the public, truth-evaluated and executed action plans.
\textbf{42/45} cells were eligible for action metrics; three
crystallization failures remain in the scheduled denominator.

Among eligible cells, \textbf{11/42} selected the true Top-1 plan, with
mean selected rank \textbf{3.31/8} and mean normalized regret
\textbf{0.297}. The joint mechanism--action outcome was more informative
than Top-1 alone: \textbf{30/42} cells had an inadequate law and a wrong
action, \textbf{11/42} had an inadequate law but a correct action,
\textbf{1/42} had an adequate law but a wrong action and \textbf{0/42}
combined an adequate law with a correct action. Law adequacy was
therefore not sufficient for action correctness, while a correct action
could occasionally occur without an adequate law.

Task heterogeneity was substantial. Electrochemical conversion reached
Top-1 in \textbf{4/15} cells, reaction safety in \textbf{4/15} and
crystallization in \textbf{3/12} eligible cells. The assay is
descriptive: only twelve task--world clusters retain complete three-arm
comparisons, and failures are concentrated in crystallization. The
bounded conclusion is a transfer boundary, not an arm-level effect.

\begin{figure*}[!t]
\centering
\includegraphics[width=\textwidth]{figures/prior-discovery/figure-6-open-action-formal.pdf}
\caption{\textbf{Unseen-action selection after exploration.} Selected rank, normalized regret, joint law--action categories and task-level means are shown for the 42 eligible cells. The three crystallization failures remain in the scheduled denominator. Every candidate plan was fully specified and executed exactly as disclosed; candidate outcomes and ranks were hidden until the terminal readout.}
\label{fig:open-action}
\end{figure*}

\section{8. Discussion}\label{discussion}

The completed evidence supports three bounded conclusions. First, priors
change search and endpoints, but their effects are context dependent: a
correct model can provide a durable advantage, an early head start or no
stable benefit, and a wrong model can still guide productive
exploration. Second, prediction improvement does not establish selective
correction. Matched evidence can trigger evidence acquisition and
numerical revision without reliable structural-law identification, and
executable summaries can lose information relative to the agent's
conditional predictions. Third, even a complete action interface does
not guarantee transfer to an unseen candidate set: only 11/42 eligible
readouts selected the true Top-1 plan, and no cell combined an adequate
law with a correct action.

These results argue against a single endpoint score for scientific
intelligence. Evidence acquisition, belief revision, law identification,
law compression and action transfer should be measured as separate
transitions. The correct interpretation is a map of conversion losses in
one controlled agent-system configuration, not a claim about all models
or all chemical settings.

\subsection{8.1 Scope and limitations}\label{scope-and-limitations}

The study uses bounded executable chemical worlds rather than universal
chemical fidelity or direct wet-laboratory validation. Five independent
worlds per task limit precision for small or heterogeneous effects. The
observation-model intervention and scope/compositionality intervention
were not executed. The matched-evidence structural block contains five
worlds, so its sign-flip result is descriptive. The open-action assay
contains 42 eligible cells, three retained crystallization failures and
twelve complete three-arm clusters; its arm summaries are not causal
estimates. Exact replay verifies the declared computational world and
plan semantics but does not remove all interface or execution burden.

\section{9. Methods}\label{methods}

\subsection{9.1 World and initial-model
construction}\label{world-and-initial-model-construction}

Each task instantiates an executable world and a participant-facing
initial model. Within a matched cluster, all arms share the world,
resource card, public contract and stochastic identity; exactly one
declared component of the initial model changes. Entity, parametric and
structural encodings are included only after separate identifiability
checks confirm comparable information volume, wording, confidence and
falsification cost.

\subsection{9.2 Campaigns, checkpoints and evaluator
scoring}\label{campaigns-checkpoints-and-evaluator-scoring}

Campaign lengths and checkpoint positions were fixed before execution. A
checkpoint records prior assessment, predictions, uncertainty, evidence
references, an executable law summary and next-experiment intent. The
evaluator owns the held-out truth queries and does not return them to
the participant. Prediction error is the mean normalized absolute error
across prespecified query--metric pairs. Final laws are executed on the
same coordinates, and blind replay compares the committed recommendation
with the observed incumbent using separate evaluator resources.

\subsection{9.3 Matched evidence and selective-correction
estimand}\label{matched-evidence-and-selective-correction-estimand}

For locus \(\ell\), the primary contrast is

\[
 C_{\ell}=(E_{\mathrm{mis,pre}}-E_{\mathrm{mis,final}})-
          (E_{\mathrm{aligned,pre}}-E_{\mathrm{aligned,final}}).
\]

The selective-correction gate requires improvement in the misspecified
arm, no unacceptable degradation in the aligned arm and a lower
confidence bound above zero. Loci are reported separately; unlike
endpoint contrasts are not pooled as one universal prior effect. Failed
scientific cells remain in the denominator and are not replaced.

\subsection{9.4 Open-action construction and
analysis}\label{open-action-construction-and-analysis}

An outcome-blind generator constructed eight complete candidate action
plans per world before participant execution. Each plan disclosed its
full operation sequence, parameters, initial-state assumptions,
measurements and terminal assay. After twelve experiments and the final
checkpoint, the participant returned a complete ranking and one selected
plan. The evaluator then calculated selected rank, Top-1, regret and the
joint law--action category. All scheduled cells remained in the
denominator.

\section{10. Completed studies and future
programme}\label{completed-studies-and-future-programme}

The following boundary is part of the paper's interpretation.

\begin{itemize}
\tightlist
\item
  \textbf{Completed --- prospective cohort:} priors change search and
  endpoints; selective correction is not established.
\item
  \textbf{Completed --- matched-evidence follow-up:} evidence
  acquisition and numerical revision can occur without exact
  structural-law recovery.
\item
  \textbf{Completed --- formal open-action assay:} unseen-action
  transfer is limited and task dependent.
\item
  \textbf{Supplementary --- exploratory cohorts and pilot:}
  developmental context, not pooled into the main denominators.
\item
  \textbf{Supplementary --- incomplete prototype, repair trajectory and
  excluded outputs:} diagnostic history; never used to replace a
  scheduled cell.
\item
  \textbf{Future --- private within-family replication:} tests stability
  on newly sealed worlds.
\item
  \textbf{Future --- context-reset artifact portability:} tests whether
  learned artifacts transfer to a new target context.
\end{itemize}

The future studies require new protocols and independent denominators.
They are not implied by the current endpoint results and are not
described as completed evidence.

\section{11. Data and code
availability}\label{data-and-code-availability}

The executable environment, prespecified protocols, analysis code,
source data and reproducible figure scripts accompany the public
release. Raw model payloads, credentials and private world instances are
excluded. Future private and context-reset studies will use newly sealed
instances and will not alter the present denominators.

\section{12. Competing interests}\label{competing-interests}

The authors declare no competing interests.

\section{Supplementary Information}\label{supplementary-information}

\subsection{S1. Early exploratory
cohorts}\label{s1.-early-exploratory-cohorts}

Before the prospective cohort, exploratory configurations were used to
test whether explicit priors altered endpoint behavior and whether the
evaluator could score prediction, law and blind replay without changing
the participant trajectory. These cohorts reached 44/45 and 75/75
terminal cell records in their respective configurations, but their task
mix, recovery rules and campaign lengths were not identical to the
prospective design. They are therefore descriptive context only. Across
the common endpoint panels, explicit information sometimes improved
outcomes over opaque identifiers, while misspecified information was not
consistently harmful. Held-out prediction error often fell, but the
wrong-prior improvement was not selectively larger and blind replay
rarely exceeded the observed incumbent.

\subsection{S2. One-world parametric
pilot}\label{s2.-one-world-parametric-pilot}

A preliminary one-world parametric pilot retained one opaque, one
aligned and one misspecified cell. All 3/3 cells, 12/12 experiments,
12/12 checkpoints, 4/4 shared truth queries and 18/18 blind replays
completed with exact replay. The misspecified agent moved outside its
supplied parameter window after contradictory evidence and reduced its
stated reliability, yet its best endpoint remained below the aligned and
opaque cells. The pilot demonstrates that model rejection and
finite-budget policy recovery are distinct; it is not a cross-task
claim.

\subsection{S3. Incomplete prototype and repair
trajectory}\label{s3.-incomplete-prototype-and-repair-trajectory}

An earlier open-action prototype retained 11/15 complete cells and did
not produce a Top-1 selection. It is not combined with the formal assay.
A separate repair trajectory completed a previously interrupted
crystallization cell but included a resource rejection and selected rank
8/8. The repair is a sensitivity result, not a replacement cell.

\subsection{S4. Operational diagnostics and excluded
outputs}\label{s4.-operational-diagnostics-and-excluded-outputs}

The development records contain schema failures, rejected operations,
recovery events, resource rejections, timing summaries and other
interface diagnostics. These quantities describe the complete agent
system and are not independent scientific samples. Implementation
counters, run identifiers, cell identifiers and raw or replaced outputs
are intentionally omitted from the reader-facing paper. Historical
outputs remain archived for auditability but cannot alter a completed
denominator.

\subsection{S5. Future studies}\label{s5.-future-studies}

Private replication will use newly sealed worlds with the same three-arm
contract and one-shot denominators. Context-reset portability will
compare raw evidence, structured evidence bundles and typed executable
laws in a new target context. Both studies require independent
preregistration and will be reported separately from the present paper.

\bibliographystyle{plainnat}
\bibliography{references}

\end{document}
